% GENERATED FILE. Edit canonical lessons, then run npm run generate:books.
% canonical-source-hash: fnv1a64:80019697497fd4ad
% canonical-lessons: MR-C545-bhavitavya, MR-C545-suvarnasandhi, MR-C545-soy, MR-C545-atkal-bandhne, MR-C545-vartavne

\chapter{Guessing the Future}
\label{ch:mr-guessing-the-future}
\hlchaptermodality{545}

\begin{chapteropening}
\textbf{By the end of this chapter:} \emph{I can talk about destiny, chances and forecasts.}

Complete the last lesson of chapter 545: I can talk about destiny, chances and forecasts.
\end{chapteropening}

\section[\texorpdfstring{\mr{भवितव्य} (bhavitavya)}{भवितव्य (bhavitavya)}]{\mr{भवितव्य} (bhavitavya) — destiny; what lies ahead}

\label{lesson:MR-C545-bhavitavya}



\begin{quote}
Before the new one: say the Marathi for to postpone, then the Marathi for to see through.
\end{quote}

\subsection*{You'll want to know: \mr{भवितव्य}}
\textbf{\mr{भवितव्य}} — \emph{bhavitavya} — \textquotedblleft{}destiny; what lies ahead\textquotedblright{}.

\begin{quote}
\mr{मुलांचं} \mr{भवितव्य} \mr{शिक्षणावर} \mr{अवलंबून} \mr{असतं}.
\end{quote}

\begin{quote}
\emph{mulāñcaṁ bhavitavya śikṣaṇāvar avlambūn astaṁ.}
\end{quote}

\begin{quote}
Children's future depends on their education.
\end{quote}

Neuter.

\begin{cousinweb}
\textbf{\mr{भवितव्य}} (\emph{bhavitavya}) is a tatsama, taken from Sanskrit \emph{bhavitavya}, \textquotedblleft{}what is to be\textquotedblright{} (Monier-Williams).
\end{cousinweb}

\begin{grammarlens}[title={What you've built}]
One more word for saying what happens if.
\end{grammarlens}

\subsection*{Your turn}
\begin{itemize}
\raggedright
  \item \emph{Say it:} \emph{bhavitavya}
  \item \emph{Say it:} \emph{bhavitavya}, once more
  \item \emph{Say it:} \emph{mulāñcaṁ bhavitavya śikṣaṇāvar avlambūn astaṁ.} — Children's future depends on their education.
  \item \emph{Recall:} say the Marathi for to postpone, then the Marathi for to see through, then say \emph{bhavitavya} again
\end{itemize}

\begin{tcolorbox}[breakable,colback=teal!4,colframe=teal!35!black,title={Before you move on}]
What does \mr{भवितव्य} mean? (\textquotedblleft{}Destiny; what lies ahead\textquotedblright{}.) Say it once more.
\end{tcolorbox}
\section[\texorpdfstring{\mr{सुवर्णसंधी} (suvarṇasandhī)}{सुवर्णसंधी (suvarṇasandhī)}]{\mr{सुवर्णसंधी} (suvarṇasandhī) — a golden chance}

\label{lesson:MR-C545-suvarnasandhi}



\begin{quote}
Before the new one: say the Marathi for to see through, then the Marathi for destiny.
\end{quote}

\subsection*{You'll want to know: \mr{सुवर्णसंधी}}
\textbf{\mr{सुवर्णसंधी}} — \emph{suvarṇasandhī} — \textquotedblleft{}a golden chance\textquotedblright{}.

\begin{quote}
\mr{ही} \mr{नोकरी} \mr{माझ्यासाठी} \mr{सुवर्णसंधी} \mr{आहे}.
\end{quote}

\begin{quote}
\emph{hī nokrī mājhyāsāṭhī suvarṇasandhī āhe.}
\end{quote}

\begin{quote}
This job is a golden chance for me.
\end{quote}

Feminine.

\begin{cousinweb}
\textbf{\mr{सुवर्णसंधी}} (\emph{suvarṇasandhī}) is a compound built from Sanskrit \emph{suvarṇa}, \textquotedblleft{}gold\textquotedblright{} (Monier-Williams), and \mr{संधी}, \textquotedblleft{}a chance\textquotedblright{}.
\end{cousinweb}

\begin{grammarlens}[title={What you've built}]
One more word for saying what happens if.
\end{grammarlens}

\subsection*{Your turn}
\begin{itemize}
\raggedright
  \item \emph{Say it:} \emph{suvarṇasandhī}
  \item \emph{Say it:} \emph{suvarṇasandhī}, once more
  \item \emph{Say it:} \emph{hī nokrī mājhyāsāṭhī suvarṇasandhī āhe.} — This job is a golden chance for me.
  \item \emph{Recall:} say the Marathi for to see through, then the Marathi for destiny, then say \emph{suvarṇasandhī} again
\end{itemize}

\begin{tcolorbox}[breakable,colback=teal!4,colframe=teal!35!black,title={Before you move on}]
What does \mr{सुवर्णसंधी} mean? (\textquotedblleft{}A golden chance\textquotedblright{}.) Say it once more.
\end{tcolorbox}
\section[\texorpdfstring{\mr{सोय} (soy)}{सोय (soy)}]{\mr{सोय} (soy) — a convenience; an arrangement}

\label{lesson:MR-C545-soy}



\begin{quote}
Before the new one: say the Marathi for destiny, then the Marathi for a golden chance.
\end{quote}

\subsection*{You'll want to know: \mr{सोय}}
\textbf{\mr{सोय}} — \emph{soy} — \textquotedblleft{}a convenience; an arrangement\textquotedblright{}.

\begin{quote}
\mr{गावात} \mr{राहण्याची} \mr{चांगली} \mr{सोय} \mr{आहे}.
\end{quote}

\begin{quote}
\emph{gāvāt rāhṇyācī cāṅglī soy āhe.}
\end{quote}

\begin{quote}
There are good arrangements for staying in the village.
\end{quote}

Feminine.

\mr{सोय} \mr{करणे}, to make arrangements.

\begin{cousinweb}
\textbf{\mr{सोय}} (\emph{soy}) is a tadbhava: CDIAL lists Marathi \emph{soy}, \textquotedblleft{}comfortableness\textquotedblright{}, under Sanskrit \emph{saukhya}, \textquotedblleft{}comfort, welfare\textquotedblright{} (CDIAL 13612).
\end{cousinweb}

\begin{grammarlens}[title={What you've built}]
One more word for saying what happens if.
\end{grammarlens}

\subsection*{Your turn}
\begin{itemize}
\raggedright
  \item \emph{Say it:} \emph{soy}
  \item \emph{Say it:} \emph{soy}, once more
  \item \emph{Say it:} \emph{gāvāt rāhṇyācī cāṅglī soy āhe.} — There are good arrangements for staying in the village.
  \item \emph{Recall:} say the Marathi for destiny, then the Marathi for a golden chance, then say \emph{soy} again
\end{itemize}

\begin{tcolorbox}[breakable,colback=teal!4,colframe=teal!35!black,title={Before you move on}]
What does \mr{सोय} mean? (\textquotedblleft{}A convenience; an arrangement\textquotedblright{}.) Say it once more.
\end{tcolorbox}
\section[\texorpdfstring{\mr{अटकळ} \mr{बांधणे} (aṭkaḷ bāndhṇe)}{अटकळ बांधणे (aṭkaḷ bāndhṇe)}]{\mr{अटकळ} \mr{बांधणे} (aṭkaḷ bāndhṇe) — to conjecture; to guess}

\label{lesson:MR-C545-atkal-bandhne}



\begin{quote}
Before the new one: say the Marathi for a golden chance, then the Marathi for a convenience.
\end{quote}

\subsection*{You'll want to know: \mr{अटकळ} \mr{बांधणे}}
\textbf{\mr{अटकळ} \mr{बांधणे}} — \emph{aṭkaḷ bāndhṇe} — \textquotedblleft{}to conjecture; to guess\textquotedblright{}.

\begin{quote}
\mr{ढग} \mr{पाहून} \mr{आजोबांनी} \mr{पावसाची} \mr{अटकळ} \mr{बांधली}.
\end{quote}

\begin{quote}
\emph{ḍhag pāhūn ājobānnī pāvsācī aṭkaḷ bāndhlī.}
\end{quote}

\begin{quote}
Seeing the clouds, Grandfather guessed it would rain.
\end{quote}

Past, with -\mr{ने} (he …): \mr{त्याने} \mr{अटकळ} \mr{बांधली}, \emph{tyāne aṭkaḷ bāndhlī}.

\begin{cousinweb}
\textbf{\mr{अटकळ} \mr{बांधणे}} (\emph{aṭkaḷ bāndhṇe}) is \mr{अटकळ}, inherited: CDIAL lists Marathi \emph{aṭkaḷ}, \textquotedblleft{}conjecture\textquotedblright{}, under \emph{*aṭṭakkalā}, \textquotedblleft{}guess\textquotedblright{} (CDIAL 183), with \mr{बांधणे}, \textquotedblleft{}to tie\textquotedblright{}.
\end{cousinweb}

\begin{grammarlens}[title={What you've built}]
One more word for saying what happens if.
\end{grammarlens}

\subsection*{Your turn}
\begin{itemize}
\raggedright
  \item \emph{Say it:} \emph{aṭkaḷ bāndhṇe}
  \item \emph{Say it:} \emph{aṭkaḷ bāndhṇe}, once more
  \item \emph{Say it:} \emph{ḍhag pāhūn ājobānnī pāvsācī aṭkaḷ bāndhlī.} — Seeing the clouds, Grandfather guessed it would rain.
  \item \emph{Recall:} say the Marathi for a golden chance, then the Marathi for a convenience, then say \emph{aṭkaḷ bāndhṇe} again
\end{itemize}

\begin{tcolorbox}[breakable,colback=teal!4,colframe=teal!35!black,title={Before you move on}]
What does \mr{अटकळ} \mr{बांधणे} mean? (\textquotedblleft{}To conjecture; to guess\textquotedblright{}.) Say it once more.
\end{tcolorbox}
\section[\texorpdfstring{\mr{वर्तवणे} (vartavṇe)}{वर्तवणे (vartavṇe)}]{\mr{वर्तवणे} (vartavṇe) — to forecast; to predict}

\label{lesson:MR-C545-vartavne}



\begin{quote}
Before the new one: say the Marathi for a convenience, then the Marathi for to conjecture.
\end{quote}

\subsection*{You'll want to know: \mr{वर्तवणे}}
\textbf{\mr{वर्तवणे}} — \emph{vartavṇe} — \textquotedblleft{}to forecast; to predict\textquotedblright{}.

\begin{quote}
\mr{हवामान} \mr{विभागाने} \mr{उद्या} \mr{पावसाचा} \mr{अंदाज} \mr{वर्तवला}.
\end{quote}

\begin{quote}
\emph{havāmān vibhāgāne udyā pāvsācā andāj vartavlā.}
\end{quote}

\begin{quote}
The weather department forecast rain for tomorrow.
\end{quote}

Past, with -\mr{ने} (he …): \mr{त्याने} \mr{अंदाज} \mr{वर्तवला}, \emph{tyāne andāj vartavlā}.

Mostly with \mr{अंदाज} or \mr{भाकीत}: \mr{अंदाज} \mr{वर्तवणे}.

\begin{cousinweb}
\textbf{\mr{वर्तवणे}} (\emph{vartavṇe}) is a learned Marathi verb, usually traced to the Sanskrit stem \emph{vart-}, \textquotedblleft{}to turn, to set going\textquotedblright{}, as in \emph{vartana}, \textquotedblleft{}setting in motion\textquotedblright{} (Monier-Williams). How it came to mean \textquotedblleft{}to forecast\textquotedblright{} is not settled.
\end{cousinweb}

\begin{grammarlens}[title={What you've built}]
One more word for saying what happens if.
\end{grammarlens}

\subsection*{Your turn}
\begin{itemize}
\raggedright
  \item \emph{Say it:} \emph{vartavṇe}
  \item \emph{Say it:} \emph{vartavṇe}, once more
  \item \emph{Say it:} \emph{havāmān vibhāgāne udyā pāvsācā andāj vartavlā.} — The weather department forecast rain for tomorrow.
  \item \emph{Recall:} say the Marathi for a convenience, then the Marathi for to conjecture, then say \emph{vartavṇe} again
\end{itemize}

\begin{tcolorbox}[breakable,colback=teal!4,colframe=teal!35!black,title={Before you move on}]
What does \mr{वर्तवणे} mean? (\textquotedblleft{}To forecast; to predict\textquotedblright{}.) Say it once more.
\end{tcolorbox}
